\section*{Objetivos Específicos}

\addcontentsline{toc}{section}{Objetivos Específicos}

\begin{enumerate}
    
    \item \textbf{Desarrollar} una interfaz web modular con autenticación y gestión de credenciales basada en Better Auth y Drizzle ORM sobre PostgreSQL/Neon, que permita al usuario parametrizar restricciones de diseño analógico en lenguaje técnico y visualizar de forma integrada el código netlist generado, la documentación funcional del circuito, el esquema resultante y los gráficos derivados de la emulación eléctrica.   
    
    \item \textbf{Entrenar e implementar} un ecosistema multiagente coordinado mediante grafos de estado en LangGraph, integrando un módulo de cálculo y un agente curador sustentado en Generación Aumentada por Recuperación (RAG) para consultar especificaciones de componentes comerciales (COTS), gobernado por un lazo de retroalimentación acotado a un presupuesto estricto de máximo 10 iteraciones de auto-corrección.

    \item \textbf{Desplegar} una arquitectura de microservicios desacoplada en contenedores Docker independientes para el cliente y el servidor stateless, configurando la pasarela de llamadas hacia APIs externas de LLMs y habilitando la ejecución segura del motor ngspice mediante PySpice en un entorno de pruebas aislado (sandbox). 
    \item \textbf{Evaluar} el desempeño global del sistema mediante un banco de pruebas de circuitos CA/CD a nivel COTS, cuantificando la tasa de convergencia en simulación (solve rate dentro de las 10 iteraciones permitidas) y la precisión de las restricciones eléctricas frente a las metas solicitadas. 
    
    
\end{enumerate}